(1) 电子获得 \(mv^2/2=eV\)，故 \(v=\sqrt{2eV/m}\)。定常电子数密度 \(n=J/(ev)\)，体电荷密度 \(\rho_q=-en=-J/v\)。使用 \(V''=-\rho_q/\varepsilon_0\)，有
\[
\boxed{V''=\frac J{\varepsilon_0}\sqrt{\frac m{2e}}\,V^{-1/2}},\quad V(0)=V'(0)=0,\ V(d)=V_d .
\]
此冷发射且临界阴极零电场条件不是任何有限电流、任意阴极温度都满足的普遍边界。

(2) 记 \(A=(J/\varepsilon_0)\sqrt{m/(2e)}\)，把微分方程乘 \(V'\)，并用阴极零场，积分得 \((V')^2=4A\sqrt V\)。取单调正根，\(V'=2\sqrt A\,V^{1/4}\)，积分 \(V^{3/4}=\frac32\sqrt A\,x\)，于是
\[
\boxed{V=V_d(x/d)^{4/3},\quad v=\sqrt{2eV_d/m}(x/d)^{2/3},\quad
J=\frac{4\varepsilon_0}{9}\sqrt{2e/m}\,\frac{V_d^{3/2}}{d^2}.}
\]
\(\rho_q=-4\varepsilon_0 V_d(x/d)^{-2/3}/(9d^2)\)。单位 \(\varepsilon_0\sqrt{e/m} V_d^{3/2}/d^2\) 是电荷/(面积\(\cdot \)时间)，幂次不是线性欧姆关系。

(3) 密度积分给
\[
Q_A=\int_0^d\rho_q\,dx=-\frac{4\varepsilon_0V_d}{3d}.
\]
局部 \(\rho_q\propto x^{-2/3}\) 虽在阴极发散，但 \(\int_0^\delta x^{-2/3}dx=3\delta^{1/3}\) 有限，故总表面面积电荷有限。电场 \(E_x=-V'=-4V_d(x/d)^{1/3}/(3d)\)，场能
\[
\boxed{U_E/A=\frac{\varepsilon_0}{2}\int_0^d E_x^2dx
=\frac{8\varepsilon_0V_d^2}{15d}},
\]
不可错误地按无空间电荷均匀电场电容套式。

(4) 粒子数通量 \(J/e\)，每粒子电势能释放并转化的动能 \(eV_d\)，因此阳极动能通量正为 \(JV_d\)，也与电场做功率 \(J\int_0^d V'dx=JV_d\) 相合。若阴极有初速或 \(V'(0)>0\)，上一问的首次积分常数不再零，\(x^{4/3}\) 与极限电流律随之改变。过短的间隙或过大的电压会引入量子隧穿、有限供电子能力、相对论电子与击穿，不能将 \(d^{-2}\) 定律无限外推。
